\documentclass[a4paper, serif, 10pt]{moderncv}

%% moderncv
% color and style
\moderncvcolor{black}
\moderncvstyle{banking}

%% page layout/margins
\usepackage[top=2.5cm, bottom=2.5cm, left=1.5cm, right=1.5cm]{geometry}

\nopagenumbers{}

%% Babel config for Italian language
% ref:
% - https://latex3.github.io/babel/guides/locale-italian.html
\usepackage[italian]{babel}

%% fontspec
\usepackage{fontspec}

\IfFontExistsTF{Linux Libertine}{
  \setmainfont[Ligatures={TeX, Common}, Numbers=OldStyle]{Linux Libertine}
}{}
\IfFontExistsTF{Linux Libertine O}{
  \setmainfont[Ligatures={TeX, Common}, Numbers=OldStyle]{Linux Libertine O}
}{}

%% CTeX
% CJK support, used to show CJK characters in the resume
%
% - heading=false: disable ctex heading
% - punct=kaiming: use kaiming punctuations styles for CJK
% - scheme=plain: use plain scheme, do not override \normalsize font size
% - space=auto: space settings for CJK characters
%
% CJK fonts are resolved per locale language, e.g. Noto Serif/Sans CJK JP
% for Japanese resumes, with fallback to the next region if not installed
%
% ref:
% - http://ctan.mirrorcatalogs.com/language/chinese/ctex/ctex.pdf
\usepackage[UTF8, heading=false, punct=kaiming, scheme=plain, space=auto]{ctex}

\IfFontExistsTF{Noto Serif CJK SC}{
  \setCJKmainfont{Noto Serif CJK SC}
}{}
\IfFontExistsTF{Noto Sans CJK SC}{
  \setCJKsansfont{Noto Sans CJK SC}
}{}

%% line spacing
\usepackage{setspace}
\setstretch{1.125}

%% URL styling - use normal text instead of monospace
\urlstyle{same}

%% Patch \httplink and \httpslink to handle full URLs
% The original moderncv macros always prepend http:// or https:// to URLs.
% This causes issues when the URL already has a protocol (e.g., https://example.com)
% resulting in malformed URLs like https://https://example.com.
% This patch checks if the URL already contains "://" and uses it directly if so.
% Note: We use \str_set:Nx (x-expansion) to expand macros like \@homepage first.
\ExplSyntaxOn
\str_new:N \g_yamlresume_url_str

\RenewDocumentCommand{\httpslink}{O{}m}{
  \str_set:Nx \g_yamlresume_url_str {#2}
  \str_if_in:NnTF \g_yamlresume_url_str {://}
    { \url{#2} }
    { \url{https://#2} }
}

\RenewDocumentCommand{\httplink}{O{}m}{
  \str_set:Nx \g_yamlresume_url_str {#2}
  \str_if_in:NnTF \g_yamlresume_url_str {://}
    { \url{#2} }
    { \url{http://#2} }
}
\ExplSyntaxOff

\name{Giulia Ferrari}{}
\title{UX Designer specializzata in ricerca utente e design system}
\phone[mobile]{+39 347 123 4567}
\email{giulia.ferrari@design.it}
\homepage{https://www.giuliaferrari.design}

\address{Via Brera 18, Milano, Lombardia, Italia, 20121}{}{}

\extrainfo{{\small \faLinkedin}\ \href{https://www.linkedin.com/in/giuliaferrari-ux}{@giuliaferrari-ux} {} {} {} • {} {} {} 
{\small \faBehance}\ \href{https://www.behance.net/giuliaferrari}{@giuliaferrari} {} {} {} • {} {} {} 
{\small \faDribbble}\ \href{https://dribbble.com/giuliaferrari}{@giuliaferrari}}

\begin{document}

\maketitle

\section{Informazioni personali}

\cvline{}{\begin{itemize}
\item UX Designer con oltre 6 anni di esperienza nella progettazione di prodotti digitali centrati sull'utente
\item Solida esperienza in ricerca qualitativa e quantitativa, usability test e analisi dei dati comportamentali
\item Progettazione di design system scalabili e accessibili, con particolare attenzione all'inclusione
\item Collaborazione con team multidisciplinari per trasformare bisogni complessi in soluzioni semplici ed efficaci
\end{itemize}}
\section{Formazione}

\cventry{set 2016–lug 2018}
        {Laurea magistrale, Design della Comunicazione, Punteggio: 110/110 con lode}
        {Politecnico di Milano}
        {\url{https://www.polimi.it}}
        {}
        {\begin{itemize}
\item Tesi sperimentale sulla progettazione di servizi pubblici digitali accessibili
\item Partecipazione a laboratori con aziende del settore fintech e healthcare
\item Approfondimento di metodologie di ricerca utente e valutazione euristica
\end{itemize}
\textbf{Corsi}: User Experience Design, Interaction Design, Metodi di ricerca qualitativa, Design System e Accessibilità, Prototipazione e Test con Utenti}

\cventry{set 2013–lug 2016}
        {Laurea triennale, Design Industriale, Punteggio: 108/110}
        {Università di Bologna}
        {\url{https://www.unibo.it}}
        {}
        {\begin{itemize}
\item Basi solide di progettazione centrata sull'utente e prototipazione rapida
\item Project work su usabilità di interfacce per elettrodomestici intelligenti
\end{itemize}
\textbf{Corsi}: Fondamenti di Design, Ergonomia e Fattori Umani, Disegno e Modellazione 3D, Storia del Design}
\section{Esperienza professionale}

\cventry{gen 2022–Presente}
        {Senior UX Designer}
        {Lumen Digital}
        {\url{https://www.lumendigital.it}}
        {}
        {\begin{itemize}
\item Guido il team UX nella progettazione di piattaforme B2B e B2C per clienti nei settori energia e retail
\item Definisco strategie di ricerca utente, conducendo interviste, usability test e analisi dei dati comportamentali
\item Progetto e mantengo design system scalabili, migliorando la coerenza dell'interfaccia e riducendo i tempi di sviluppo
\item Collaboro con product manager, sviluppatori e stakeholder per tradurre bisogni complessi in soluzioni semplici
\end{itemize}
\textbf{Parole chiave}: Ricerca utente, Design system, Usability test, Leadership UX}

\cventry{mar 2019–dic 2021}
        {UX Designer}
        {Studio Forma}
        {\url{https://www.studioforma.it}}
        {}
        {\begin{itemize}
\item Progettato flussi e interfacce per app mobile e siti web in ambito turismo, cultura e servizi pubblici
\item Condotto workshop di co-design con utenti e stakeholder per validare concept e priorità di prodotto
\item Realizzato wireframe, prototipi interattivi e specifiche di design per team di sviluppo agile
\item Contribuito alla creazione di linee guida di accessibilità e design inclusivo
\end{itemize}
\textbf{Parole chiave}: Prototipazione, Co-design, Accessibilità, Mobile UX}

\cventry{set 2018–feb 2019}
        {Junior UX Designer}
        {Agenzia Pixel}
        {\url{https://www.agenziapixel.it}}
        {}
        {\begin{itemize}
\item Supportato il team di design nella ricerca utente e nella mappatura dei customer journey
\item Creato wireframe e mockup per campagne digitali e landing page
\item Partecipato a revisioni di usabilità e alla documentazione dei risultati
\end{itemize}
\textbf{Parole chiave}: Customer journey, Wireframe, Mockup}
\section{Lingue}

\cvline{Italiano}{Madrelingua o bilingue \hfill \textbf{Parole chiave}: Madrelingua}
\cvline{Inglese}{Competenza professionale completa \hfill \textbf{Parole chiave}: IELTS 8.0, Scrittura UX}
\cvline{Francese}{Competenza professionale limitata \hfill \textbf{Parole chiave}: DELF B1}
\section{Competenze}

\cvline{Ricerca UX}{Esperto \hfill \textbf{Parole chiave}: Interviste utente, Usability test, Analisi dati, Journey mapping, Personas}
\cvline{UI Design}{Avanzato \hfill \textbf{Parole chiave}: Figma, Sketch, Adobe XD, Prototipazione, Accessibilità}
\cvline{Design System}{Avanzato \hfill \textbf{Parole chiave}: Componenti, Token di design, Documentazione, Governance}
\cvline{Collaborazione}{Avanzato \hfill \textbf{Parole chiave}: Workshop, Agile, Stakeholder, Mentoring}
\cvline{Analisi}{Intermedio \hfill \textbf{Parole chiave}: Google Analytics, Hotjar, Maze, A/B testing}
\section{Premi}

\cventry{nov 2023}
        {Associazione Design Italia}
        {Premio Design Accessibile}
        {}
        {}
        {Riconoscimento per il progetto di riqualificazione dell'esperienza utente di un servizio pubblico digitale, con particolare attenzione all'accessibilità e all'inclusione.}
\section{Certificazioni}

\cventry{giu 2021}
        {Google}
        {Google UX Design Professional Certificate}
        {\url{https://www.coursera.org/professional-certificates/google-ux-design}}
        {}
        {}

\cventry{feb 2023}
        {Nielsen Norman Group}
        {Certificazione Nielsen Norman Group - UX Research}
        {\url{https://www.nngroup.com/ux-certification/}}
        {}
        {}
\section{Pubblicazioni}

\cventry{apr 2022}
        {Progettare servizi pubblici digitali accessibili}
        {Interaction Design Magazine Italia}
        {\url{https://www.interactiondesign.it/articoli/servizi-pubblici-accessibili}}
        {}
        {\begin{itemize}
\item Articolo su metodologie di ricerca e co-design applicate alla Pubblica Amministrazione
\item Analisi di casi studio italiani e linee guida per test di usabilità inclusivi
\end{itemize}}
\section{Progetti}

\cventry{mar 2022–set 2022}
        {Progetto end-to-end di redesign dell'app mobile per una banca digitale.}
        {Ridisegno dell'app di home banking Aurora}
        {\url{https://www.lumendigital.it/progetti/aurora}}
        {}
        {\begin{itemize}
\item Condotto ricerca utente con 24 clienti e analisi dei dati di utilizzo
\item Progettato nuovi flussi di pagamento e gestione delle finanze personali
\item Ridotto del 35\% le chiamate al supporto per le operazioni più comuni
\end{itemize}
\textbf{Parole chiave}: Mobile banking, Ricerca utente, Redesign, Design system}

\cventry{gen 2020–giu 2021}
        {Sistema di design modulare per portali e app di promozione culturale.}
        {Design system per il turismo culturale}
        {\url{https://www.studioforma.it/progetti/turismo-culturale}}
        {}
        {\begin{itemize}
\item Definito libreria di componenti accessibili e linee guida editoriali
\item Allineato il design system alle WCAG 2.1 AA
\item Formato il team interno su utilizzo e manutenzione dei componenti
\end{itemize}
\textbf{Parole chiave}: Design system, Accessibilità, Cultura, WCAG}
\section{Interessi}

\cvline{Fotografia}{Paesaggio, Ritratto}
\cvline{Trekking}{Alpi, Sentieri}
\cvline{Musica}{Pianoforte, Jazz}
\section{Volontariato}

\cventry{set 2021–Presente}
        {Mentore UX}
        {Designers Italia}
        {\url{https://designers.italia.it}}
        {}
        {\begin{itemize}
\item Affiancato team di giovani designer nella progettazione di servizi pubblici digitali
\item Organizzato sessioni di formazione su ricerca utente e accessibilità
\item Contribuito a linee guida per test di usabilità con utenti fragili
\end{itemize}}
    
\end{document}